% !TeX program = pdflatex
\documentclass[11pt,a4paper]{article}
\usepackage{../../recursos/latex/estilo-notas}

\begin{document}
\cabecalhoaula{10}{Máquinas de Vetores de Suporte (SVM)}{AME §8.2 e §4.6; ISLP cap.~9}

%==============================================================================
\section*{Objetivos da aula}
%==============================================================================
Ao final desta aula o aluno deve ser capaz de:
\begin{itemize}[leftmargin=1.4cm]
  \item definir hiperplano separador e o \textbf{classificador de margem máxima};
  \item identificar os \textbf{vetores de suporte} e o papel da margem;
  \item generalizar para dados não separáveis com o \textbf{classificador de vetor
        de suporte} (\emph{soft margin}) e o hiperparâmetro $C$;
  \item explicar o \textbf{truque do \emph{kernel}} a partir da forma dual;
  \item usar \emph{kernels} (polinomial, gaussiano/RBF) para fronteiras não
        lineares e entender $\gamma$;
  \item reconhecer a SVM como minimização da \emph{perda \emph{hinge}} regularizada
        e sua relação com a regressão logística; tratar o caso multiclasse.
\end{itemize}

\begin{emsala}
Mudança de perspectiva em relação à Aula~08: lá modelávamos $\Prob(Y\mid\x)$; a SVM
ataca a fronteira \emph{geometricamente}, buscando a separação ``mais folgada''. É
um dos métodos mais elegantes do curso --- liga geometria, otimização e os RKHS da
Aula~04. Pergunta de abertura: \emph{``se há infinitos hiperplanos que separam as
classes, qual é o melhor?''}
\end{emsala}

%==============================================================================
\section{Hiperplanos e o classificador de margem máxima}
%==============================================================================
Codifique as classes como $y\in\{-1,+1\}$. Um \textbf{hiperplano} em $\R^p$ é o
conjunto $\{\x: f(\x)=0\}$ com $f(\x)=\beta_0+\bbeta^\top\x$. Ele separa o espaço
em dois lados conforme o sinal de $f$. Dizemos que os dados são \emph{linearmente
separáveis} se existe $f$ tal que
\begin{equation}\label{eq:sep}
  y_i\,f(\X_i) > 0 \qquad \text{para todo } i.
\end{equation}
A magnitude $\abs{f(\x)}$ (com $\norm{\bbeta}=1$) é a \emph{distância} de $\x$ ao
hiperplano.

\begin{ideia}
Quando há \emph{vários} hiperplanos separadores, a SVM escolhe o de \textbf{maior
margem}: a maior distância mínima aos pontos. A intuição: deixar a maior ``folga''
possível entre as classes tende a generalizar melhor.
\end{ideia}

\begin{definicao}[Classificador de margem máxima]
Procura-se o hiperplano que maximiza a margem $M$:
\[
  \max_{\beta_0,\bbeta}\ M \quad\text{s.a.}\quad
  \sum_{j=1}^p \beta_j^2 = 1,\qquad
  y_i\,f(\X_i)\ge M \ \ \forall i.
\]
A restrição $\norm{\bbeta}=1$ torna os hiperplanos comparáveis (assim $y_if(\X_i)$
é a distância com sinal). As observações que \emph{atingem} a margem
($y_if(\X_i)=M$) são os \textbf{vetores de suporte}.
\end{definicao}

\begin{atencao}
Apenas os vetores de suporte determinam o hiperplano: mover um ponto longe da
margem não muda nada. Por outro lado, a formulação \emph{exige} separabilidade
perfeita e é \textbf{muito sensível} a pequenas mudanças nos dados --- uma única
observação pode alterar drasticamente a solução. Precisamos relaxá-la.
\end{atencao}

%==============================================================================
\section{O classificador de vetor de suporte (\emph{soft margin})}
%==============================================================================
Para lidar com dados não separáveis (e ganhar robustez), permitimos que algumas
observações violem a margem, introduzindo \emph{folgas} $e_i\ge 0$:
\begin{equation}\label{eq:soft}
  \max_{\beta_0,\bbeta,\,\bm{e}}\ M \quad\text{s.a.}\quad
  \norm{\bbeta}=1,\quad
  y_i\,f(\X_i)\ge M(1-e_i)\ \ \forall i,\quad
  e_i\ge 0,\ \ \sum_{i=1}^n e_i \le C.
\end{equation}
A folga $e_i$ mede o quanto a $i$-ésima observação viola a margem: $e_i>0$ significa
dentro da margem; $e_i>1$ significa do \emph{lado errado} do hiperplano (erro de
classificação).

\begin{ideia}
$C\ge 0$ é o \textbf{tuning parameter} (orçamento de violações), escolhido por
validação cruzada (Aula~03):
\begin{itemize}[leftmargin=1.2cm,nosep]
  \item $C$ \emph{grande}: tolera muitas violações $\Rightarrow$ margem larga, mais
        vetores de suporte, modelo mais \emph{rígido} (alto viés, baixa variância);
  \item $C$ \emph{pequeno}: poucas violações $\Rightarrow$ margem estreita, modelo
        mais \emph{flexível} (baixo viés, alta variância).
\end{itemize}
É o balanço viés--variância da Aula~01 de novo. (Atenção: a convenção de $C$ no
\texttt{scikit-learn} é \emph{invertida} --- lá $C$ grande penaliza mais as
violações.)
\end{ideia}

%==============================================================================
\section{O truque do \emph{kernel}}
%==============================================================================
A chave que torna a SVM poderosa: a solução depende dos dados \emph{apenas através
de produtos internos}. Pode-se mostrar que a função ótima se escreve
\begin{equation}\label{eq:dual}
  f(\x) = \beta_0 + \sum_{k=1}^n \alpha_k\,\inner{\x}{\X_k},
  \qquad \inner{\x}{\X_k}=\textstyle\sum_{j=1}^p x_j x_{k,j},
\end{equation}
e que o cálculo dos $\alpha_k$ também só requer os produtos internos
$\inner{\X_i}{\X_k}$ entre as observações. Crucialmente, $\alpha_k=0$ para todos os
pontos que \emph{não} são vetores de suporte; logo
\[
  f(\x) = \beta_0 + \sum_{k\in S}\alpha_k\,\inner{\x}{\X_k},
\]
onde $S$ é o conjunto de vetores de suporte.

\begin{ideia}
Como tudo depende só de produtos internos, podemos \textbf{substituir}
$\inner{\x}{\X_k}$ por um \emph{kernel} $K(\x,\X_k)$ --- que corresponde a um produto
interno num espaço de \emph{features} de dimensão muito maior (ou infinita), sem
\emph{nunca} computar esse espaço explicitamente. Isso é o \textbf{truque do
\emph{kernel}}: fronteiras não lineares no espaço original a custo de fronteiras
lineares num espaço ampliado.
\end{ideia}

\noindent \emph{Kernels} de Mercer comuns (cf.\ Aula~04, §4.6):
\begin{center}\small
\renewcommand{\arraystretch}{1.4}
\begin{tabular}{@{}ll@{}}
\toprule
\textbf{Kernel} & \textbf{$K(\x,\x')$} \\
\midrule
Linear      & $\inner{\x}{\x'}$ \\
Polinomial  & $(1+\inner{\x}{\x'})^{p}$ \\
Gaussiano (RBF) & $\exp\!\big(-\gamma\,\norm{\x-\x'}^2\big)$ \\
\bottomrule
\end{tabular}
\end{center}
No kernel gaussiano, $\gamma>0$ é um \emph{tuning parameter}: $\gamma$ grande
$\Rightarrow$ influência local, fronteiras muito flexíveis (risco de superajuste);
$\gamma$ pequeno $\Rightarrow$ fronteiras suaves. Escolhem-se $C$ e $\gamma$ juntos
por validação cruzada.

%==============================================================================
\section{SVM como perda \emph{hinge} regularizada}
%==============================================================================
Há uma leitura que conecta a SVM a tudo o que vimos. Dado um kernel de Mercer $K$
com RKHS $\mathcal{H}_K$ (Aula~04), o classificador da SVM minimiza
\begin{equation}\label{eq:hinge}
  \min_{g\in\mathcal{H}_K}\ \sum_{k=1}^n \underbrace{\big(1-y_k\,g(\X_k)\big)_+}_{\text{perda \emph{hinge}}}
  + \lambda\,\norm{g}_{\mathcal{H}_K}^2 ,
\end{equation}
onde $(u)_+=\max(u,0)$. Ou seja, a SVM é \emph{ajuste de perda + regularização}
(Aula~02), com a \textbf{perda \emph{hinge}} $L(y,g)=(1-yg)_+$.

\begin{observacao}
A perda \emph{hinge} é \textbf{zero} quando $y\,g(\x)\ge 1$, isto é, quando o ponto
está do lado certo e \emph{folgadamente} longe da margem. Só os pontos próximos ou
violando a margem contribuem --- e são exatamente os \emph{vetores de suporte}. No
caso separável, \eqref{eq:hinge} reencontra a margem máxima ($1/\norm{g}$).
\end{observacao}

\subsection{Relação com a regressão logística}
A SVM (\emph{hinge}) e a regressão logística (perda logística
$\log(1+e^{-yg})$) são ambas ``perda de margem $+$ regularização $\ell_2$''. As
perdas são parecidas, mas:
\begin{itemize}[leftmargin=1.4cm,nosep]
  \item a \emph{hinge} é exatamente zero longe da margem $\Rightarrow$ solução
        \emph{esparsa} em vetores de suporte; a logística nunca é zero;
  \item a logística entrega \textbf{probabilidades} calibradas (Aula~09); a SVM
        entrega só a fronteira (probabilidades exigem pós-processamento, ex.\
        \emph{Platt scaling}).
\end{itemize}
Regra prática: classes bem separadas $\to$ SVM costuma ir muito bem; quando se quer
$\Prob(Y\mid\x)$, prefira a logística.

%==============================================================================
\section{Mais de duas classes}
%==============================================================================
A SVM é intrinsecamente binária. Para $\abs{\mathcal{C}}>2$:
\begin{description}[leftmargin=2.8cm,style=nextline]
  \item[Um-contra-um (OvO)] treina uma SVM para cada par de classes
        ($\binom{K}{2}$ classificadores) e classifica por votação. Padrão do
        \texttt{scikit-learn}.
  \item[Um-contra-todos (OvA)] treina $K$ SVMs (cada classe vs.\ o resto) e escolhe
        a de maior $f(\x)$.
\end{description}

%==============================================================================
\section{Prática em Python}
%==============================================================================
\begin{lstlisting}
from sklearn.svm import SVC
from sklearn.preprocessing import StandardScaler
from sklearn.pipeline import make_pipeline
from sklearn.model_selection import GridSearchCV

# SVM com kernel RBF; padronizar e ESSENCIAL (metodo baseado em distancia)
pipe = make_pipeline(StandardScaler(), SVC(kernel="rbf"))
grade = {"svc__C":     [0.1, 1, 10, 100],
         "svc__gamma": [0.001, 0.01, 0.1, 1]}
busca = GridSearchCV(pipe, grade, cv=5, scoring="accuracy")
busca.fit(X_tr, y_tr)
print("melhores parametros:", busca.best_params_)

# Para probabilidades: SVC(probability=True) (usa Platt scaling, mais lento)
\end{lstlisting}
Os notebooks \texttt{Aula prática 10} e \texttt{Exemplo - SVM em dados artificiais}
ilustram o efeito de $C$,
$\gamma$ e da escolha do kernel sobre a fronteira de decisão.

\begin{atencao}
SVM com kernel RBF é \textbf{sensível à escala} das covariáveis (depende de
distâncias) --- \emph{sempre} padronize dentro do \emph{pipeline} (Aula~07).
\end{atencao}

%==============================================================================
\section*{Resumo}
%==============================================================================
\begin{itemize}[leftmargin=1.4cm]
  \item \textbf{Margem máxima}: hiperplano separador de maior folga; definido pelos
        \emph{vetores de suporte}; exige separabilidade e é frágil.
  \item \textbf{\emph{Soft margin}} \eqref{eq:soft}: folgas $e_i$ e orçamento $C$
        (tuning) controlam o balanço viés--variância.
  \item \textbf{Truque do kernel}: a solução só depende de produtos internos
        \eqref{eq:dual}; trocá-los por $K$ dá fronteiras não lineares (polinomial,
        RBF com $\gamma$).
  \item SVM $=$ \textbf{perda \emph{hinge}} $+$ regularização \eqref{eq:hinge}
        (RKHS, Aula~04); esparsa em vetores de suporte.
  \item Relação com a logística (perda de margem $+$ $\ell_2$); SVM não dá
        probabilidades de imediato.
  \item Multiclasse por OvO/OvA; \emph{padronize} sempre.
\end{itemize}

\paragraph{Leitura recomendada.} [AME] §8.2 (SVM) e §4.6 (RKHS e truque do kernel);
[ISLP] Capítulo~9 (§9.1 margem máxima, §9.2 classificador de vetor de suporte, §9.3
SVM com kernels, §9.4 multiclasse, §9.5 relação com a logística).

\end{document}
